\section*{第19章：案例研究——Featherweight Java}
\addcontentsline{toc}{section}{第19章：案例研究——Featherweight Java}

\begin{taplsolutionentrybox}
\phantomsection
\label{sol:translator:19.4.3}
\textbf{练习19.4.3}

FJ 原本把对象值直接写成
\(\mathsf{new}\ C(\overline v)\)。这样的值没有独立身份；若要修改字段，
必须像引用演算那样另设存储。下面给出一种完整扩展。

在项和值中加入位置 \(l\) 和赋值：
\[
 t::=\cdots\mid l\mid t.f:=t
 \qquad
 v::=l
\]
存储 \(\mu\) 把位置映射到已完成初始化的对象
\(\mathsf{new}\ C(\overline v)\)。求值判断改写为
\((t,\mu)\trans(t',\mu')\)。对象创建先按从左到右的次序求值全部实参，再分配
新位置：
\[
\inferrule*[right=\textsc{E-NewStore}]
 {l\notin\operatorname{dom}(\mu)}
 {(\mathsf{new}\ C(\overline v),\mu)
  \trans
  (l,\mu[l\mapsto\mathsf{new}\ C(\overline v)])}
\]
字段读取改为从存储取值：
\[
\inferrule*[right=\textsc{E-FieldStore}]
 {\mu(l)=\mathsf{new}\ C(\overline v)\\
  \mathit{fields}(C)=\overline C\ \overline f}
 {(l.f_i,\mu)\trans(v_i,\mu)}
\]
方法调用仍使用 \(\mathit{mbody}\)，但现在以接收者位置替换
\texttt{this}：
\[
\inferrule*[right=\textsc{E-InvkStore}]
 {\mu(l)=\mathsf{new}\ C(\overline v)\\
  \mathit{mbody}(m,C)=(\overline x,t_0)}
 {(l.m(\overline u),\mu)
  \trans
  ([\overline x\mapsto\overline u,\mathsf{this}\mapsto l]t_0,\mu)}
\]
类型转换必须查询接收者位置中保存对象的运行期类：
\[
\inferrule*[right=\textsc{E-CastStore}]
 {\mu(l)=\mathsf{new}\ C(\overline v)\\C\subtype D}
 {((D)l,\mu)\trans(l,\mu)}
\]
若 \(C\nsubtype D\)，这项转换没有可用规则，仍按 FJ 的约定卡住。缺少这条规则
会使对象一经分配成位置后，所有本来可以成功的类型转换都无法继续。

赋值先求值接收者，再求值右侧；两者均为值后，按照以下规则更新对应字段：
\[
\inferrule*[right=\textsc{E-Assign}]
 {\mu(l)=\mathsf{new}\ C(v_1,\ldots,v_i,\ldots,v_n)\\
  \mathit{fields}(C)=C_1f_1,\ldots,C_nf_n}
 {(l.f_i:=v,\mu)
  \trans
  (v,\mu[l\mapsto\mathsf{new}\ C(v_1,\ldots,v,\ldots,v_n)])}
\]
规则结论右侧配置的第一分量是 \(v\)，表示赋值表达式本身的求值结果就是刚写入
的值；第二分量则是完成字段更新后的存储。这与 Java 赋值表达式的行为一致。若
希望赋值只是一条命令，也可加入 \texttt{Unit} 并返回 \texttt{unit}。

所有同余规则都要改成携带存储。可以用求值上下文一次写全：在
\cref{def:19.5.3}的上下文语法中加入
\(E.f:=t\) 与 \(l.f:=E\)，并使用
\[
\inferrule*[right=\textsc{E-ContextStore}]
 {(t,\mu)\trans(t',\mu')}
 {(E[t],\mu)\trans(E[t'],\mu')}
\]
这里的上下文仍规定从左到右求值，因此字段访问、方法调用、对象创建、类型转换
和赋值都会把每一步产生的新存储继续传给外围项。

类型层面加入存储类型 \(\Sigma\)，其中 \(\Sigma(l)=C\) 表示位置 \(l\) 存放
类 \(C\) 的对象。位置和赋值规则为
\[
\inferrule*[right=\textsc{T-Loc}]
 {\Sigma(l)=C}{\Gamma\mid\Sigma\vdash l:C}
\]
\[
\inferrule*[right=\textsc{T-Assign}]
 {\Gamma\mid\Sigma\vdash t_0:C_0\\
  \mathit{fields}(C_0)=\overline C\ \overline f\\
  \Gamma\mid\Sigma\vdash t_1:D\\
  D\subtype C_i}
 {\Gamma\mid\Sigma\vdash t_0.f_i:=t_1:C_i}
\]
还需定义 \(\Sigma\vdash\mu\)：二者定义域相同，并且每个位置中对象的全部字段值
都符合该类声明的字段类型。保持性定理随后与\chapref{13}相同：一步求值可能
扩展存储和存储类型，但已有位置的类型不变；字段更新保持存储良类型。

\taplsolutionbacklink{ex:19.4.3}{返回练习19.4.3}
\end{taplsolutionentrybox}

\begin{taplsolutionentrybox}
\phantomsection
\label{sol:translator:19.4.4}
\textbf{练习19.4.4}

加入 Java 中用于抛出异常的 \texttt{throw}，以及用于捕获异常的
\texttt{try}--\texttt{catch} 构造：
\[
 t::=\cdots
 \mid \mathsf{throw}\ t
 \mid \mathsf{try}\ t\ \mathsf{catch}\ (C\ x)\ t
\]
并约定可抛出的对象都是某个根类 \texttt{Throwable} 的子类。异常本身仍是普通
对象值。当被抛出的项已经求值为对象值 \(v\) 时，\(\mathsf{throw}\ v\) 表示
异常已经产生；接下来会跳过外层尚未完成的计算，向外寻找类型匹配的
\texttt{catch}。

先求值 \texttt{throw} 的参数与 \texttt{try} 的受保护项：
\[
\inferrule*[right=\textsc{E-Throw}]
 {t\trans t'}{\mathsf{throw}\ t\trans\mathsf{throw}\ t'}
\qquad
\inferrule*[right=\textsc{E-Try}]
 {t\trans t'}
 {\mathsf{try}\ t\ \mathsf{catch}\ (C\ x)\ h
  \trans
  \mathsf{try}\ t'\ \mathsf{catch}\ (C\ x)\ h}
\]
受保护项正常得到值时，处理器被丢弃：
\[
 \mathsf{try}\ v\ \mathsf{catch}\ (C\ x)\ h\trans v
\]
若得到 \(\mathsf{throw}\ v\)，令 \(D\) 为 \(v\) 的运行时类。若
\(D\subtype C\)，执行处理器并以 \(v\) 替换 \(x\)；否则继续向外传播：
\[
\inferrule*[right=\textsc{E-Catch}]
 {v=\mathsf{new}\ D(\overline v)\\D\subtype C}
 {\mathsf{try}\ (\mathsf{throw}\ v)\ \mathsf{catch}\ (C\ x)\ h
  \trans[x\mapsto v]h}
\]
\[
\inferrule*[right=\textsc{E-Miss}]
 {v=\mathsf{new}\ D(\overline v)\\D\nsubtype C}
 {\mathsf{try}\ (\mathsf{throw}\ v)\ \mathsf{catch}\ (C\ x)\ h
  \trans\mathsf{throw}\ v}
\]
此外，对字段访问、方法调用、对象创建、类型转换及其参数位置，都要加入异常
传播规则；一旦当前求值位置产生 \(\mathsf{throw}\ v\)，就停止求值同一构造中
尚未处理的部分，并把异常传到外层。

\(\mathsf{throw}\ t\) 一旦执行，就不会向当前计算返回普通值，而会中断正常
求值并向外寻找处理器；因此，它不可能产生一个类型不符合外层预期的正常结果。
这使它可以出现在任何预期类型的位置，例如类型声明为 \(C\) 但总是抛出异常的
方法体。相应的类型规则让 \(\mathsf{throw}\ t\) 可以取得任意预期类型：
\[
\inferrule*[right=\textsc{T-Throw}]
 {\Gamma\vdash t:D\\D\subtype\mathsf{Throwable}}
 {\Gamma\vdash\mathsf{throw}\ t:C}
\]
前提只要求被抛出的对象具有 \texttt{Throwable} 的子类型；结论中的 \(C\) 与
异常对象的类型 \(D\) 无关，可以由 \(\mathsf{throw}\ t\) 所在位置的类型要求
决定。

处理器的两个正常出口必须具有相容类型。采用 FJ 的算法式风格，可要求处理器
结果是受保护项结果类型的子类型，或反过来，再选择较大的类型；采用带包容规则的
声明式系统，则可先把两个分支都提升到共同超类型 \(T\)：
\[
\inferrule*[right=\textsc{T-Try}]
 {\Gamma\vdash t:T\\
  \Gamma,x:C\vdash h:T\\
  C\subtype\mathsf{Throwable}}
 {\Gamma\vdash
  \mathsf{try}\ t\ \mathsf{catch}\ (C\ x)\ h:T}
\]
这样，异常传播不会改变外围项期望的类型；进展性则把最终结果扩展为三类：普通
值、未处理的 \(\mathsf{throw}\ v\)，或失败的类型转换。

\taplsolutionbacklink{ex:19.4.4}{返回练习19.4.4}
\end{taplsolutionentrybox}

\begin{taplsolutionentrybox}
\phantomsection
\label{sol:translator:19.4.5}
\textbf{练习19.4.5}

可以加入包容规则，但类型转换不能直接读取一个已经经过包容提升的类型。先记
原书的算法判断为 \(\Gamma\vdash_{\min}t:S\)：它为项计算唯一的最小类型
\(S\)。再定义声明式判断 \(\Gamma\vdash_{\mathrm d}t:T\)，并加入
\[
\inferrule*[right=\textsc{T-Sub}]
 {\Gamma\vdash_{\mathrm d}t:S\\S\subtype T}
 {\Gamma\vdash_{\mathrm d}t:T}
\]
对方法调用、对象创建和方法体检查，可以删去原规则中单列的子类型前提，直接
借助 T-Sub 把实参或方法体提升到声明所要求的类型。例如：
\[
\inferrule*[right=\textsc{T-InvkD}]
 {\Gamma\vdash_{\mathrm d}t_0:C_0\\
  \mathit{mtype}(m,C_0)=\overline D\to C\\
  \Gamma\vdash_{\mathrm d}\overline t:\overline D}
 {\Gamma\vdash_{\mathrm d}t_0.m(\overline t):C}
\]
\[
\inferrule*[right=\textsc{T-NewD}]
 {\mathit{fields}(C)=\overline D\ \overline f\\
  \Gamma\vdash_{\mathrm d}\overline t:\overline D}
 {\Gamma\vdash_{\mathrm d}\mathsf{new}\ C(\overline t):C}
\]
若实参的最小类型为 \(C_i\subtype D_i\)，先由其声明式规则取得 \(C_i\)，再用
T-Sub 提升为 \(D_i\)，即可满足这些规则。

向上、向下和荒谬转换必须依据同一个确定的源类型分类；否则，T-Sub 可以先改变
被转换项的类型，再让同一转换落入另一条规则。因此，三条转换规则都以算法检查器
算出的最小类型为准：
\[
\inferrule*[right=\textsc{T-UCastD}]
 {\Gamma\vdash_{\min}t_0:D\\D\subtype C}
 {\Gamma\vdash_{\mathrm d}(C)t_0:C}
\]
\[
\inferrule*[right=\textsc{T-DCastD}]
 {\Gamma\vdash_{\min}t_0:D\\C\subtype D\\C\ne D}
 {\Gamma\vdash_{\mathrm d}(C)t_0:C}
\]
第三条荒谬转换规则不再另列；它也使用前提
\(\Gamma\vdash_{\min}t_0:D\)，并要求 \(C\) 与 \(D\) 互不为子类型。

实现仍先运行原书的算法检查器求最小类型；声明式推导只说明此后允许把结果用于
任何超类型位置。若算法为项 \(t\) 算出最小类型 \(S\)，声明式系统便可将
\(t\) 赋予 \(S\) 或 \(S\) 的任意超类型 \(T\)；反过来，声明式系统能够赋予
\(t\) 的每个类型 \(T\) 都满足 \(S\subtype T\)。这样，T-Sub 只会增加项可以
取得的超类型，不会改变类型转换的判断。例如，虽然
\(\mathsf{new}\ B():B\) 可以通过 T-Sub 提升为 \texttt{Object}，转换
\((A)\mathsf{new}\ B()\) 仍依据最小类型 \(B\) 检查，因此依旧是荒谬转换，
不会被误判为从 \texttt{Object} 到 \(A\) 的普通向下转换。

\taplsolutionbacklink{ex:19.4.5}{返回练习19.4.5}
\end{taplsolutionentrybox}

\begin{taplsolutionentrybox}
\phantomsection
\label{sol:translator:19.5.5}
\textbf{练习19.5.5}

下面给出一个可执行的 FJ 实现。完整程序由类表、静态检查和求值三部分组成；
完整源码位于
\path{code/book-snippets/src/chapter19.rs}。这道工程题的实现较长，书中不再
逐页展开 Rust 源码，而是交代数据模型、关键检查、求值顺序和端到端例子，
使读者可以先理解实现结构，需要时再到工程中逐函数核对。

\textbf{一、数据模型与类表。}
项的数据类型直接对应\cref{fig:19-1}中的五种形式：变量、字段访问、方法调用、
对象创建和类型转换；类声明则保存类名、超类、字段、构造函数与方法。类表是
从类名到类声明的映射，并提供四项基本查询：
\begin{description}
\item[\(\mathit{subtype}(C,D)\)] 从 \(C\) 沿超类链向上查找，判断是否能到达
\(D\)。
\item[\(\mathit{fields}(C)\)] 先取得继承字段，再追加 \(C\) 自己声明的字段。
\item[\(\mathit{mtype}(m,C)\)] 沿继承链查找方法 \(m\)，返回形参类型序列与
结果类型。
\item[\(\mathit{mbody}(m,C)\)] 按相同次序查找方法，返回形参名序列与方法体。
\end{description}
载入类表时还要检查类名和成员名不重复、引用的类已经定义、继承图无环、子类
不遮蔽继承字段，以及覆盖方法的形参与结果类型和原方法完全相同。

\textbf{二、类声明检查。}
构造函数的形参必须依次列出继承字段与当前类字段；传给 \texttt{super} 的实参
必须恰好是继承字段；随后每个当前类字段都由同名形参赋给 \texttt{this.f}。
检查方法时，在包含各形参和 \texttt{this} 的环境中求出方法体类型，确认它是
声明结果类型的子类型，再检查方法覆盖是否满足上述约束。

\textbf{三、项的类型检查。}
类型检查函数按\cref{fig:19-4}递归处理项。变量从环境取类型；字段访问先求出
接收者类型，再从该类的完整字段序列中查找字段；方法调用先查询方法签名，然后
逐项检查实参与形参；对象创建则用完整字段序列检查构造实参。类型转换依据源类与
目标类的子类型关系分为向上转换、向下转换和荒谬转换；实现记录荒谬转换警告，
但仍按原 FJ 规则返回目标类型。

\textbf{四、单步求值。}
求值函数按\cref{fig:19-3}从左到右寻找下一项待求值的子项。接收者已经成为
\(\mathsf{new}\ C(\overline v)\) 后，字段访问按字段在完整字段序列中的位置
取值；方法调用查询方法体，并同时完成替换
\(\{x_i\mapsto v_i,\mathsf{this}\mapsto v_0\}\)。这里必须使用同时替换，
否则某个实参中恰好出现另一形参名时，后续替换会错误地再次改写它。运行期对象的
类是目标类的子类型时，类型转换可以删去；否则该项卡住。多步解释器反复执行一次
单步求值，直到得到对象值或再无规则可用。

\textbf{五、端到端例子。}
正文中的 \texttt{Pair.setfst} 调用依次求值为
\[
\begin{aligned}
 &\mathsf{new}\ \mathsf{Pair}(\mathsf{new}\ A(),\mathsf{new}\ B()).
   \mathsf{setfst}(\mathsf{new}\ B())\\
 \trans{}&\mathsf{new}\ \mathsf{Pair}
   (\mathsf{new}\ B(),
    (\mathsf{new}\ \mathsf{Pair}(\mathsf{new}\ A(),\mathsf{new}\ B())).\mathsf{snd})\\
 \trans{}&\mathsf{new}\ \mathsf{Pair}(\mathsf{new}\ B(),\mathsf{new}\ B())
\end{aligned}
\]
第一步把实参替换给 \texttt{newfst}，并把接收者替换给 \texttt{this}；第二步
再从接收者对象中投影出 \texttt{snd} 字段。这条路径同时经过方法查询、同步替换
和字段投影，因而可以作为整份实现的最小端到端检查。

配套测试还覆盖合法与失败的类型转换、荒谬转换警告、构造函数和方法参数错误、
非法覆盖、继承环，以及重复或保留名称。完整源码参加 \texttt{make rust-check}
的编译和测试；书中只保留理解实现所需的结构、算法与完整求值路径。

\taplsolutionbacklink{ex:19.5.5}{返回练习19.5.5}
\end{taplsolutionentrybox}

\begin{taplsolutionentrybox}
\phantomsection
\label{sol:translator:19.5.6}
\textbf{练习19.5.6}

GJ 风格扩展的核心，是让类、方法和类型都可以带有受上界约束的类型变量。下面
说明如何扩展\cref{ex:19.5.5}的实现，并列出新增机制的关键代码。

\textbf{语法。}
类型从类名扩展为类型变量或带实参的类类型：
\[
 T::=X\mid C\langle\overline T\rangle
\]
类与方法声明分别增加类型形参
\(\langle\overline X\subtype\overline N\rangle\)，对象创建、方法调用和类型
转换也携带相应的类型实参。上界 \(N\) 可以提到先前声明的类型变量，从而表达
F-有界约束，例如 \(X\subtype\mathit{Comparable}\langle X\rangle\)。

\taplrustsupportflowchunkfirstcompact{ch19-gj-syntax}{1}{20}
\taplrustsupportflowchunk{ch19-gj-syntax}{21}{38}

\textbf{环境。}
在变量环境 \(\Gamma\) 之外增加类型变量环境 \(\Delta\)，记录每个类型变量的
上界。先定义类型良构判断 \(\Delta\vdash T\ \mathsf{OK}\)：类类型的实参数量
必须正确，且把实际类型代入声明上界后，每个实参都必须是对应上界的子类型。

类或方法的形式类型变量换成实际类型实参时，递归使用下面的类型替换函数。

\taplrustsupportflowchunkfirstcompact{ch19-gj-type-substitution}{1}{15}

\textbf{子类型。}
增加变量规则
\[
 \inferrule{X\subtype N\in\Delta}{\Delta\vdash X\subtype N}
\]
类继承规则需要先把类声明中的形式类型变量替换为实际类型参数。例如，若
\[
 \mathsf{class}\ C\langle\overline X\subtype\overline N\rangle
 \ \mathsf{extends}\ D\langle\overline U\rangle\{\ldots\},
\]
则
\[
 C\langle\overline T\rangle
 \subtype
 [\overline X\mapsto\overline T]
 D\langle\overline U\rangle
\]
类型参数本身保持不变；GJ 的类类型并不会仅因实参之间有子类型关系就自动形成
协变关系。

下面的实现同时处理类型变量的上界、沿继承链进行的子类型查询，以及类类型实参
是否满足声明上界的检查。

\taplrustsupportflowchunkfirstcompact{ch19-gj-subtyping}{1}{25}
\taplrustsupportflowchunk{ch19-gj-subtyping}{26}{55}
\taplrustsupportflowchunk{ch19-gj-subtyping}{56}{84}

\textbf{字段与方法查询。}
\(\mathit{fields}\)、\(\mathit{mtype}\) 和 \(\mathit{mbody}\) 沿继承链查找
时，都必须累积并应用类型替换。泛型方法调用还要先检查显式类型实参满足方法的
上界，再把它们代入形参类型、返回类型与方法体中的类型标注。它们和上一题的查询
框架相同，主要差别就是多了这一步替换，故不再重复列出整段实现。

\textbf{类型检查与求值。}
类型判断改写为 \(\Delta;\Gamma\vdash t:T\)。对象创建和方法调用先检查类型
实参，再按替换后的字段或形参类型检查项实参。运行时仍可采用 FJ 的求值规则：
类型参数不参与字段选择和方法体替换，可以在执行前擦除；但证明擦除正确，需要
先证明类型替换保持良构与赋型，再证明项替换保持赋型。

这里采用完整的擦除执行路径。建立泛型类表时先检查类型形参、上界、继承、字段、
方法签名和方法体；字段或方法的接收者是类型变量时，沿 \(\Delta\) 中的上界查询
成员。执行程序前，把泛型类声明、字段、方法签名和方法体全部擦除为 FJ 类表，
为每个类生成 FJ 的标准构造函数。若类型实参代入后，某个字段或方法的结果类型
比其擦除后声明的结果类型更具体，擦除过程会在成员访问外插入必要的 FJ 类型转换；
例如，\texttt{Box<A>.get()} 的声明在擦除后返回 \texttt{Object}，用于需要
\texttt{A} 的位置时便插入到 \texttt{A} 的转换。最后以同样方式擦除待求值项，
并交给上一题的 FJ 解释器实际执行。

下面的代码给出擦除的关键步骤。\texttt{erase\_type} 递归把类型变量替换成其上界
的擦除；该函数只有这一条递归和类名保持不变两种情形，故不再单独列出。
\texttt{erase\_typed\_term} 先保留源项的静态类型，再在字段访问或方法调用
的擦除结果过于宽泛时插入 FJ 转换。它调用的字段类型与方法结果类型查询只负责沿
擦除后的继承链查找声明，前文已经说明其算法，故不再重复列出。对象创建和原有转换
只需递归擦除类型参数。

\taplrustsupportflowchunkfirstcompact{ch19-gj-erasure-term}{1}{30}
\taplrustsupportflowchunk{ch19-gj-erasure-term}{31}{58}

\Needspace{14\baselineskip}
\textbf{验证顺序。}
先测试无泛型 FJ 程序在扩展后行为不变；再测试单一泛型容器、带界类型变量、泛型
方法、继承时改变类型实参、递归上界，以及违反上界的拒绝案例。还要拒绝重定义
\texttt{Object} 和改变类型形参上界的非法方法覆盖，并用泛型返回值处于更具体类型
上下文的程序检验转换插入。最后比较擦除前后的求值结果。这样既覆盖 GJ 扩展新增
的静态机制，也验证它没有改变 FJ 原有程序的可观察行为。

完整实现中的 \texttt{GenericTable::check\_classes} 检查类和方法声明；
\texttt{fields}、\texttt{method}、\texttt{type\_of} 与
\texttt{check\_arguments} 完成成员查询和类型检查；\texttt{erase\_method}、
\texttt{erase\_table} 与 \texttt{eval\_erased} 把整个程序交给 FJ 解释器执行。
测试覆盖无泛型程序回归、泛型容器和方法、上界成员查询、改变类型实参的继承、
递归上界、上界违反、非法覆盖，以及擦除后需要插入转换的方法体和链式成员调用。

\taplsolutionbacklink{ex:19.5.6}{返回练习19.5.6}
\end{taplsolutionentrybox}
